\section{结论}
\label{sec:conclusion}

我们识别出动作 token VLA 特有的结构性安全缺口：其主干通常是未对齐的 base LLM，而其动作词汇经过数千步行为克隆训练后，驻留在自然语言拒绝方向不再覆盖的残差流子空间中。\textsc{RDT}+ 通过从外部对齐兄弟 LLM 移植拒绝方向，并经由覆盖全部 7 个动作 token 的生成阶段感知双阶段策略注入，在其应处的位置恢复拒绝信号。方向来源对照实验证实，移植的拒绝方向的特定几何性质——而非仅仅是隐层状态扰动的幅度——是有效成分。\textsc{RDT}+ 无需训练，增加不超过 5\% 的延迟，可跨动作攻击类型泛化，并可与提示层防御组合使用。我们将 \textsc{RDT}+ 视为将文本对齐视为跨非文本输出通路\emph{可重用资产}的一步，并呼吁学界将 base LLM 主干和词表覆写明确视为需要分析的结构性安全决策。
